\chapter{测试章节1}

\section{测试节标题1}

\begin{definition}{上道1}{shangdao1}
	content
\end{definition}

\begin{axiom}{上道2}{shangdao2}
	content
\end{axiom}

\begin{hypothesis}{上道3}{shangdao3}
	content
\end{hypothesis}

\begin{theorem}{上道4}{shangdao4}
	content
\end{theorem}

\begin{lemma}{上道5}{shangdao5}
	content
\end{lemma}

\begin{corollary}{上道6}{shangdao6}
	content
\end{corollary}

\begin{metatheorem}{上道7}{shangdao7}
	content
\end{metatheorem}

\begin{criteria}{上道8}{shangdao8}
	content
\end{criteria}

\begin{problem}{上道9}{shangdao9}
	content
\end{problem}

\autoref{sample-alarm}应该是警示而不是定理,\autoref{shangdao2}应该是公理而不是定理.


\autoref{sample-algorithm}

\begin{example}{上道10}{shangdao10}
	content
\end{example}

\begin{remark}{上道11}{shangdao11}
	content
\end{remark}

\begin{proof}
	content
\end{proof}

\begin{answer}
	content
\end{answer}

\begin{sketch}{}{}
	content
\end{sketch}

\begin{hint}
	提示环境现在支持自然换行、分段和跨页。

	这是提示中的第二段。
\end{hint}

\begin{algorithm}{示例算法}{sample-algorithm}
	输入数据，执行若干步骤，然后返回结果。
\end{algorithm}

\begin{convention}{符号约定}{sample-convention}
	本文统一采用这一记号约定。
\end{convention}

\begin{alarm}{重要警示}{sample-alarm}
	这里放置容易出错或必须特别注意的内容。
\end{alarm}

\begin{code}[Python 示例]
def square(x):
    return x * x
\end{code}
